\section{Related Literature}\label{sec:related}

This paper relates to three strands of literature: acquisitions and innovation, entrepreneurial finance and startup policy, and local versus national policy evaluation.

The closest connection is to the literature on acquisitions of innovative entrants and their effects on innovation and competition. \citet{cunningham2021killer} show that acquisitions can eliminate future competition by shutting down overlapping innovation projects. Subsequent theoretical work has broadened that perspective. \citet{letina2024killer} study how merger policy affects innovation strategies more generally, while \citet{dijk2024direction} and \citet{dijk2025strategic} analyze how acquisition prospects distort startup innovation and R\&D incentives. Our paper builds on these insights but asks a different question: how should entry subsidies be designed when buyout exit and merger review jointly shape project choice and the social value of successful entry?

The paper also connects to the literature on value capture in startup markets. \citet{arora2024missing} emphasize that startups often rely on external buyers or downstream partners to realize returns. That perspective is central here: buyout exit is not a secondary outcome but a determinant of private entrepreneurial incentives before the product market is reached.

A second related literature studies entrepreneurial finance and public startup support. \citet{lerner2020vcrole}, \citet{lerner2022governmentincentives}, and \citet{bai2021public} discuss venture capital, public entrepreneurial finance, and government support for startups. Empirical work evaluates instruments such as R\&D tax credits, investor tax credits, accelerators, multilevel startup support, and broader startup policy packages \citep{fazio2020rdtax,denes2023investortax,gonzalezuribe2018accelerators,lanahan2015multilevel,manaresi2021startupact}. Our paper is complementary to that literature: rather than estimating treatment effects, it develops a theory of why the effectiveness of startup subsidies depends on the acquisition environment and why stronger private entry incentives need not imply stronger grounds for subsidy.

The third relevant literature concerns geography and the divergence between local and national policy evaluation. \citet{chen2010buylocal} and \citet{chatterji2014clusters} motivate attention to local entrepreneurial ecosystems. \citet{kline2013peopleplaces} provide a foundational framework for understanding why local and national welfare effects of place-based policy can differ. Our mechanism is different. We do not model strategic subsidy competition across jurisdictions. Instead, we study how buyout exit changes the incidence of gains and losses from startup policy, thereby generating a local--national incidence wedge even when the underlying private equilibrium is the same.

Relative to these literatures, the paper's novelty is not simply to add acquisitions to a startup-subsidy model. It combines three elements that are usually studied separately. First, it embeds endogenous merger attempt into both private incentives and welfare accounting, yielding a primitive benchmark subsidy formula driven by product-market primitives and the merger-attempt condition. Second, it makes entrepreneurial direction a choice variable, so merger permissiveness changes not only how much entrepreneurship occurs but also what type occurs; a uniform subsidy cannot correct that composition distortion. Third, it shows that the local--national wedge is not only an incidence comparison for a fixed project: because buyout exit can trigger project switching, the economically relevant wedge can itself change with the exit environment.
